\section{Method}

\subsection{Preliminaries}

\paragraph{Notation.}
Let $\mathbf{x} = (x_1, x_2, \ldots, x_N)$ denote a sequence of $N$ tokens, and let $\mathtt{M}$ denote the special mask token.
We denote the parameter of language model as $\theta$.
Notably, both autoregressive and diffusion-based language models share the same Transformer architecture, they differ only in their input representations and training objectives.

\paragraph{Autoregressive Language Models.}
Autoregressive (AR) language models is trained to minimize the negative log-likelihood of the sequence:
\begin{equation}
    \mathcal{L}_{\text{AR}} = -\mathbb{E}_{\mathbf{x} \sim \mathcal{D}} \left[ \sum_{i=1}^{N} \log p_{\theta}(x_i \mid x_{<i}) \right].
\end{equation}

\paragraph{Masked Diffusion Language Models.}
Masked diffusion language models (MDLMs)~\citep{llada} define a discrete diffusion process over token sequences.
Let $t\sim\mathcal{U}(0,1)$ denote a noise level, and let
$\mathbf{x}^t$ denote the corrupted sequence obtained by independently
masking each token $x_i$ with probability $t$.
Formally, for each position $i$:
\begin{equation}
    x_i^{t} =
    \begin{cases}
        \mathtt{M} & \text{with probability } t, \\
        x_i & \text{otherwise}.
    \end{cases}
\end{equation}
The model is trained to predict the original tokens from the corrupted sequence:
\begin{equation}
\begin{aligned}
&\mathcal{L}_{\mathrm{MDLM}}
=
-\mathbb{E}_{t,\mathbf{x},\mathbf{x}^{t}}
\left[
\frac{1}{t}
\sum_{i=1}^{N}
\mathbb{I}\!\left[x_i^{t}=\mathtt{M}\right]
\ell_{\theta,i}^{\mathrm{MDLM}}(t)
\right], \\
&\ell_{\theta,i}^{\mathrm{MDLM}}(t)
=
\log p_{\theta}
\left(
x_i
\,\middle|\,
\mathbf{x}^{t}
\right).
\end{aligned}
\label{eq:mdlm_loss}
\end{equation}


\paragraph{Block Diffusion Language Models.}
Block diffusion language models (BDLMs)~\citep{bd3lms} extend MDLMs by
partitioning the sequence into contiguous blocks of size $B$.
The sequence is divided into $N_b=N/B$ blocks,
$\mathbf{x}=(\mathbf{b}_1,\ldots,\mathbf{b}_{N_b})$.
During training, each block $r\in\{1,\ldots,N_b\}$ is assigned an
independently sampled noise level $t_r\sim\mathcal{U}(0,1)$.
Tokens within $\mathbf{b}_r$ are then masked independently with probability
$t_r$, while the causal structure across blocks is maintained.
For each token $x_i$, let
$j=\lceil i/B\rceil$
denote its block index in the following objective:

\begin{equation}
\begin{aligned}
&\mathcal{L}_{\mathrm{BDLM}}
=
-\mathbb{E}_{\boldsymbol{t},\mathbf{x},\mathbf{x}^{\boldsymbol{t}}}
\left[
\sum_{i=1}^{N}
\frac{
\mathbb{I}\!\left[x_i^{\boldsymbol{t}}=\mathtt{M}\right]
}{
t_j
}
\ell_{\theta,i}^{\mathrm{BDLM}}(t_j)
\right], \\
&
\ell_{\theta,i}^{\mathrm{BDLM}}(t_j)
=
\log p_{\theta}
\left(
x_i
\,\middle|\,
\mathbf{b}_{<j},
\mathbf{b}_{j}^{t_j}
\right).
\end{aligned}
\label{eq:bdlm_loss}
\end{equation}

where $\mathbf{b}_{<j}$ denotes the clean preceding blocks in $\mathbf{x}$, and $\mathbf{b}_{j}^{t}$ denotes the noisy version of the $j$-th block in $\mathbf{x}^{t}$. The specialized attention mask used for efficient BDLM training is illustrated in Figure~\ref{fig:original_pattern}.

\subsection{\name{}: \NAME{}}

\paragraph{Training formulation}
Given a sequence $\mathbf{x}$, we sample $K$ independent noise levels $t_1, t_2, \ldots, t_K \in [0, 1]$ and generate $K$ corresponding corrupted sequences $\mathbf{x}^{t_1}, \mathbf{x}^{t_2}, \ldots, \mathbf{x}^{t_K}$.
The training objective of \name{} is:
\begin{equation}
\begin{aligned}
    &\mathcal{L}_{\text{\name{}}} = -\mathbb{E}_{\{t_k\}_{k=1}^K, \mathbf{x}, \{\mathbf{x}^{t_k}\}_{k=1}^K} \Bigg[\\
    &{1\over K}\sum_{k=1}^K {1\over t_k} \sum_{i=1}^N [x_i^{t_k} = \mathtt{M}] \log p_{\theta}(x_i \mid \mathbf{b}_{<block(i)}, \mathbf{b}_{block(i)}^{t_k}) \Bigg].
\end{aligned}
\end{equation}

\paragraph{Training formulation}
Given a clean sequence $\mathbf{x}$, \name{} shares this clean sequence as the common context and generates $K$ corrupted sequences for denoising training.
For each view $k\in\{1,\ldots,K\}$, a noise level $t_k\sim\mathcal{U}(0,1)$ is independently sampled, and the corresponding corrupted sequence $\mathbf{x}^{t_k}$ is generated by masking each token with probability $t_k$.
The training objective averages the denoising losses over these $K$ noised views:
\begin{equation}
\begin{aligned}
&\mathcal{L}_{\mathrm{MMT}}
=
\mathbb{E}
\left[
\frac{1}{K}
\sum_{k=1}^{K}
\mathcal{J}_{k}(\mathbf{x})
\right], \\
&\mathcal{J}_{k}(\mathbf{x})
=
-\frac{1}{t_k}
\sum_{i=1}^{N}
\mathbb{I}\!\left[x_i^{t_k}=\mathtt{M}\right]
\ell_{\theta,i}^{\mathrm{MMT}}(t_k), \\
&\ell_{\theta,i}^{\mathrm{MMT}}(t_k)
=
\log p_{\theta}
\left(
x_i
\,\middle|\,
\mathbf{b}_{<j},
\mathbf{b}_{j}^{t_k}
\right).
\end{aligned}
\label{eq:mmt_loss}
\end{equation}
where the expectation is taken over $\mathbf{x}$, the independently sampled noise levels $\{t_k\}_{k=1}^{K}$, and the corresponding corrupted sequences $\{\mathbf{x}^{t_k}\}_{k=1}^{K}$.
Here, $j$ denotes the block index of token $x_i$, $\mathbf{b}_{<j}$ denotes the clean preceding blocks shared by all noised views, and $\mathbf{b}_{j}^{t_k}$ denotes the noisy version of the $j$-th block in the $k$-th corrupted sequence.

where the expectation is taken over $\mathbf{x}$, the block-wise noise levels
$\{\boldsymbol{t}_k\}_{k=1}^{K}$, and the corresponding corrupted views
$\{\mathbf{x}^{\boldsymbol{t}_k}\}_{k=1}^{K}$.
The block notation follows Eq.~\eqref{eq:bdlm_loss}, while
$t_{k,j}$ denotes the noise level assigned to block $j$ in the $k$-th corrupted view.

\paragraph{Context-sharing attention mechanism}
To enable efficient parallel training of all $K$ perturbations while maintaining correct token visibility, we construct an attention mask $\mathbf{M} \in \{0, 1\}^{(K+1)N \times (K+1)N}$ over token indices.
We use $k = 0$ for the clean sequence $\mathbf{x}$ and $k \in \{1, \ldots, K\}$ to index the $K$ perturbations $\mathbf{x}^{t_k}$.
Let $(k, i)$ denote the $i$-th token of the $k$-th sequence.
The attention mask entry $M_{(k,i), (k',i')}$ determines whether token $(k,i)$ can attend to token $(k',i')$:
\begin{equation}
    M_{(k,i), (k',i')} =
    \begin{cases}
        1, & \text{if } k' = 0 \text{ and } \lceil i' / B \rceil < \lceil i / B \rceil, \\
        1, & \text{if } k' = k \text{ and } \lceil i' / B \rceil = \lceil i / B \rceil, \\
        0, & \text{otherwise}.
    \end{cases}
\label{eq:attn_mask}
\end{equation}
This mask ensures that each perturbation can attend to: (1) preceding blocks from the shared clean sequence, and (2) its own noisy block at the current position.
The efficient training mask of \name{} is illustrated in Figure~\ref{fig:method}.

\paragraph{Complementary masking via Context Sharing.}
We adopt the complementary masking strategy~\citep{fastdllmv2}.
Instead of independently sampling $K$ mask patterns, we sample $K/2$ pairs of complementary mask patterns.
Crucially, unlike prior implementations that double the context computation cost, our context-sharing paradigm allows us to process complementary masks against a single clean context representation.
 The training objective of \name{} becomes:
\begin{equation}
    \mathcal{L}_{\text{\name{}}} = -\mathbb{E}_{\{t_j\}_{j=1}^{K/2}, \mathbf{x}, \{\mathbf{x}^{t_j}\}_{j=1}^{K/2}} \Bigg[
    \frac{1}{K} \sum_{j=1}^{K/2} \mathcal{L}_{\text{pair}}^{(j)}\Bigg].
\end{equation}
Each pair of complementary masks consist an inverse mask patterns $m$ and $\bar m= 1-m$, generating two noisy sequece $\mathbf{x}^{t_j}$ and $\bar{\mathbf{x}}^{t_j}$, where
\begin{equation}
    \bar{x}_i^{t_j} =
    \begin{cases}
        x_i & \text{if } x_i^{t_j} = \mathtt{M}, \\
        \mathtt{M} & \text{otherwise}.
    \end{cases}
\end{equation}
The paired loss is then defined as
\begin{equation}
\begin{aligned}
    \mathcal{L}_{\text{pair}}^{(j)} = & \sum_{i=1}^N [x_i^{t_j} = \mathtt{M}] \log p_{\theta}(x_i \mid \mathbf{b}_{<block(i)}, \mathbf{b}_{block(i)}^{t_j}) \\
    + & \sum_{i=1}^N [\bar{x}_i^{t_j} = \mathtt{M}] \log p_{\theta}(x_i \mid \mathbf{b}_{<block(i)}, \bar{\mathbf{b}}_{block(i)}^{t_j}) \Bigg],
\end{aligned}
\end{equation}
where $\bar{\mathbf{b}}^{t_j}_{block(i)}$ is the $i$-th noisy blocks of $\bar{\mathbf{x}}^{t_j}$. We eliminate the normalization factor $1/t_j$ following \citet{fastdllmv2}, as the complementary masks jointly cover all $N$ tokens.

\paragraph{Training algorithm.}
We summarize the complete training procedure of \name{} in Algorithm~\ref{alg:mmt}.

(Working in Progress...)

\begin{algorithm}[t]
\caption{\name{} Training}
\label{alg:mmt}
\begin{algorithmic}[1]
\REQUIRE Model $p_\theta$, dataset $\mathcal{D}$, number of mask pairs $K/2$, block size $B$
\FOR{each training iteration}
    \STATE Sample minibatch $\{\mathbf{x}\}$ from $\mathcal{D}$
    \STATE \textcolor{gray}{\textit{// Generate K complementary mask pairs}}
    \FOR{$j = 1$ to $K/2$}
        \STATE Sample $t_j \sim \text{Uniform}(0, 1)$
        \STATE Sample $\mathbf{m}_j \in \{0,1\}^N$, $m_{j,i} \sim \text{Bernoulli}(t_j)$
        \STATE $\mathbf{x}^{t_j} \gets \mathbf{x} \odot (1 - \mathbf{m}_j) + \mathtt{M} \odot \mathbf{m}_j$
        \STATE $\bar{\mathbf{x}}^{t_j} \gets \mathbf{x} \odot \mathbf{m}_j + \mathtt{M} \odot (1 - \mathbf{m}_j)$
    \ENDFOR
    \STATE \textcolor{gray}{\textit{// Construct input with shared clean context}}
    \STATE $\mathbf{X} \gets \text{Concat}(\mathbf{x}, \mathbf{x}^{t_1}, \bar{\mathbf{x}}^{t_1}, \ldots, \mathbf{x}^{t_{K/2}}, \bar{\mathbf{x}}^{t_{K/2}})$
    \STATE Construct attention mask $\mathbf{M}$ per Eq.~(\ref{eq:attn_mask})
    \STATE \textcolor{gray}{\textit{// Forward and backward pass}}
    \STATE $\mathcal{L} \gets \mathcal{L}_{\text{\name{}}}(p_\theta(\mathbf{X}; \text{attn\_mask}=\mathbf{M}))$
    \STATE Update $\theta$ via gradient descent on $\mathcal{L}_{\text{\name{}}}$
\ENDFOR
\end{algorithmic}
\end{algorithm}

\subsection{Computational Complexity Analysis}

We analyze the computational complexity of different language modeling paradigms.
Let $N$ denote the sequence length, $P$ the total model parameters, $L$ the number of transformer layers, and $d$ the hidden dimension.

\paragraph{Per-token training cost.}
Following \citet{chen-etal-2025-cost}, the FLOPs required to process one token with context length $n$ is:
\begin{equation}
    C_{\text{token}}(n) = \underbrace{2P}_{\text{non-attention}} + \underbrace{4n L d}_{\text{attention}}.
\end{equation}
Here $n$ is the number of visible tokens in the masked attention mechanism.

\paragraph{Autoregressive Transformers.}
For causal language models generating a sequence of length $N$, the total FLOPs is:
\begin{equation}
\begin{aligned}
    C_{\text{AR}} &= \sum_{t=1}^{N} C_{\text{token}}(t) = 2PN + \sum_{t=1}^{N} 4t L d \\
    &= \underbrace{2PN}_{\text{non-attention}} + \underbrace{2N(N+1) L d}_{\text{attention}}.
\end{aligned}
\end{equation}

\paragraph{Block Diffusion Language Models.}
Block diffusion models partition the sequence into $N_b = N / B$ blocks.
At block $j$ of the clean sequence $\mathbf{x}$, the model attends to all $(j-1)B$ preceding tokens plus the current block of size $B$.
At block $j$ of the noisy sequence $\mathbf{x}^{t}$, the model attends to all $(j-1)B$ preceding clean tokens plus the current noisy block of size $B$.
That is, the $j$-th block of each sequence attends to $jB$ tokens.
The total FLOPs is
\begin{equation}
\begin{aligned}
    C_{\text{BDLM}} &= 2 \sum_{i=1}^{N} C_{\text{token}}((\lceil i/B \rceil)B) \\
    &= 2 \sum_{j=1}^{N_b} B \cdot C_{\text{token}}(jB) \\
    &= 2 (2P N_b B + 2 B^2 N_b(N_b+1) L d) \\
    &= 2\times \left(\underbrace{2PN}_{\text{non-attention}} + \underbrace{2N(N+B) L d}_{\text{attention}}\right).
\end{aligned}
\end{equation}

\paragraph{\name{} (Ours).}
Our context-sharing mechanism processes the clean context once and shares it across $K$ perturbations.
Similar to the above, we process $1+K$ sequences in total, where the $j$-th block of each sequence attends to $jB$ tokens.
The total FLOPs is
\begin{equation}
\begin{aligned}
    C_{\text{\name{}}} &= (1+K) \sum_{j=1}^{N_b} B \cdot C_{\text{token}}(jB) \\
    &= (1+K) \times \left(\underbrace{2PN}_{\text{non-attention}} + \underbrace{2N(N+B) L d}_{\text{attention}}\right).
\end{aligned}
\end{equation}
The amortized cost per perturbation is:
\begin{equation}
    \frac{C_{\text{\name{}}}}{K} = (1+1/K) \times \left(\underbrace{2PN}_{\text{non-attention}} + \underbrace{2N(N+B) L d}_{\text{attention}}\right).
\end{equation}

In typical training settings, the sequence length $N$ is at least 2048 or 4096, while the block size $B$ is usually set to small values such as 32 or 64~\citep{fastdllmv2}.
Since $B \ll N$, we have $N(N+B) \approx N^2$ and $N(N+1) \approx N^2$.
Table~\ref{tab:flops_comparison} summarizes the average training FLOPs per noise perturbation under this approximation.
As shown, \name{} significantly improves the FLOPs efficiency compared to block diffusion, which can be viewed as a special case of our method with $K=1$.

\begin{table}[h]
\centering
\caption{Average FLOPs per sequence/perturbation. For \name{}, the cost is amortized over $K$ perturbations, and $(1+1/K) \approx 1$ when $K$ is reasonably large.}
\label{tab:flops_comparison}
\scalebox{0.9}{
\begin{tabular}{l|ccc}
\toprule
 & AR & BDLM & \name{} (Ours) \\
\midrule
Non-attention & $2PN$ & $4PN$ & $(1+1/K) \cdot 2PN$ \\
Attention & $2N^2 Ld$ & $4N^2 Ld$ & $(1+1/K) \cdot 2N^2 Ld$ \\
\bottomrule
\end{tabular}
}
\end{table}
